\documentclass[a4paper]{article}

\usepackage{graphicx}
\usepackage{amsmath,amssymb}
\usepackage{minted}
\usepackage{multirow}
\usepackage{float}
\usepackage{todonotes}
\usepackage{forest}
\usepackage{xstring}
\usepackage{xcolor}
\usepackage[most]{tcolorbox}
\usepackage{xparse}
\usepackage{natbib}

\usetikzlibrary{graphs,graphdrawing}
\usegdlibrary{layered}

% for external graphviz
\input{/home/donkarlo/Dropbox/repo/nd_ai_project/src/nd_ai/knoweldge_representation/language/formal/computer/programming/tex/latex/code_clips/external_graphviz.tex}

% for tags
\input{/home/donkarlo/Dropbox/repo/nd_ai_project/src/nd_ai/knoweldge_representation/language/formal/computer/programming/tex/latex/parts/control_sequence/kind/semtag/semtag.tex}

% for links
\input{/home/donkarlo/Dropbox/repo/nd_ai_project/src/nd_ai/knoweldge_representation/language/formal/computer/programming/tex/latex/code_clips/seehere.tex}

\title{Gloabl Workspace theory}
\date{}

\begin{document}
    \maketitle
    \clickurlfooter{https://en.wikipedia.org/wiki/Global_workspace_theory} argues that the brain contains many specialized processes or modules that operate in parallel, much of which is unconscious. The global workspace is a functional hub of broadcast and integration that allows information to be disseminated across modules. As such GWT can be classified as a functionalist theory of consciousness. When sensory input, memories, or internal representations receive attention, they enter the global workspace and become accessible to various cognitive processes. As elements compete for attention, those that succeed gain entry to the global workspace, allowing their information to be distributed and coordinated throughout the whole cognitive system. GWT resembles the concept of working memory and is proposed to correspond to a 'momentarily active, subjectively experienced' event in working memory. It facilitates top-down control of attention, working memory, planning, and problem-solving through this information sharing.
    %\bibliographystyle{plainnat}
    %\bibliography{/home/donkarlo/Dropbox/repo/nd_shared_research/bibliography/bibliography.bib}
\end{document}